\chapter*{Abstract}
% \addcontentsline{toc}{chapter}{Abstract}
\label{chap:abstract}

Accurate traffic flow prediction, particularly short-term traffic speed forecasting, is essential for Intelligent Transportation Systems (ITS), enabling applications such as route planning, congestion management, and emission reduction. Traditional approaches often rely on centralized architectures that require continuous transmission of raw sensor data to a central server, raising concerns about communication overhead, potential privacy and governance concerns, and scalability. Federated Learning can mitigate these limitations by collaboratively training models across edge devices without sharing raw data directly. However, standard federated algorithms such as FedAvg can suffer from degraded performance in traffic sensor networks because traffic data are highly non-Independent and Identically Distributed (non-IID) across heterogeneous road segments.

This thesis proposes a communication-efficient Hierarchical Federated Learning framework for traffic flow prediction in ITS. The framework combines four components: Dynamic Time Warping (DTW)-based clustering to group traffic sensors by temporal pattern similarity; two-tier hierarchical aggregation to balance intra-cluster specialization with DTW-distance-weighted inter-cluster knowledge sharing; quality-weighted aggregation to prioritize client updates with lower local training loss; and adaptive reliability-based client selection to reduce communication cost under dynamic client availability.

The framework is evaluated on the METR-LA dataset, containing 207 traffic sensors, using a lightweight GRU Sequence-to-Sequence model and a seven-configuration cumulative ablation study. The final configuration achieves an overall Mean Absolute Error (MAE) of 3.676~mph, corresponding to a 4.1\% improvement over the global FedAvg baseline. Compared with the most directly comparable federated baseline, CMBA-FL, the proposed framework reduces MAE by 13.5\% while using 29.1\% less communication, requiring 25.8~GB compared with 36.4~GB. Adaptive client selection reduces communication by 12.2\% with only a 0.4\% accuracy trade-off. Cross-dataset validation on PEMS-BAY, containing 325 sensors, achieves an MAE of 2.052~mph without dataset-specific architectural tuning, providing initial evidence for the generalizability of the proposed approach.

The results show that temporal-pattern-aware clustering is the largest individual contributor to accuracy improvement in the ablation study. Overall, the proposed framework offers a favorable communication--accuracy trade-off compared with existing federated traffic forecasting methods, while avoiding centralized raw-data collection and explicit graph-structure requirements.

\vspace{1cm}
\noindent\textbf{Keywords:} Federated Learning, Hierarchical Aggregation, Traffic Flow Prediction, Dynamic Time Warping, Communication Efficiency, Non-IID Data, Intelligent Transportation Systems, GRU Seq2Seq